% !TeX program = xelatex
\documentclass[aspectratio=169,11pt]{beamer}
\usepackage[UTF8,fontset=none]{ctex}
\setCJKmainfont{Microsoft YaHei}
\setCJKsansfont{Microsoft YaHei}
\setsansfont{Arial}
\usepackage{amsmath,amssymb,booktabs,array}
\definecolor{ink}{HTML}{17314D}
\definecolor{muted}{HTML}{586575}
\definecolor{accent}{HTML}{146A86}
\setbeamercolor{normal text}{fg=ink,bg=white}
\setbeamercolor{frametitle}{fg=ink}
\setbeamercolor{structure}{fg=accent}
\setbeamerfont{frametitle}{size=\Large,series=\bfseries}
\setbeamertemplate{navigation symbols}{}
\setbeamertemplate{itemize items}[circle]
\setbeamertemplate{footline}{\hfill{\color{muted}\scriptsize\insertframenumber}\hspace{6mm}\vspace{3mm}}
\setbeamersize{text margin left=9mm,text margin right=9mm}
\setlength{\parskip}{5pt}
\newcommand{\source}[1]{{\vfill\scriptsize\color{muted}#1\par}}
\newcommand{\key}[1]{{\color{accent}\bfseries#1}}
% Generated by update_tables.py; do not edit numbers manually.

\newcommand{\PhaseOneTable}{
\begin{tabular}{@{}lrrr@{}}
\toprule
被保护任务 & 独立 & 无保护AGA后 & 保护AGA后 \\
\midrule
MNIST & 99.69 & 99.32 & 9.80 \\
DTD & 79.41 & 68.78 & 2.13 \\
SVHN & 97.46 & 91.77 & 6.70 \\
\bottomrule
\end{tabular}
}

\newcommand{\MainTable}{
\begin{tabular}{@{}lrrrr@{}}
\toprule
方法 & 独立MNIST & 直接融合 & \shortstack{AGA处理\\保护专家} & \shortstack{AGA处理\\全部专家} \\
\midrule
无保护 & 99.69 & 99.37 / 92.50 & 99.30 / 92.78 & 99.32 / 91.77 \\
偏置扰动$^*$ & 99.69 & 9.80 / 6.70 & 9.80 / 6.70 & 9.80 / 6.70 \\
补偿 O/B$^\dagger$ & 99.71 & 78.77 / 62.94 & 79.53 / 65.82 & 77.85 / 61.42 \\
补偿 A/B$^\dagger$ & 99.65 & 91.07 / 77.46 & 85.42 / 75.50 & 79.20 / 69.84 \\
PaRaMS$^*$ & 99.69 & 8.92 / 9.15 & 99.30 / 92.79 & 99.32 / 91.77 \\
\bottomrule
\end{tabular}
}

\newcommand{\UtilityTable}{
\begin{tabular}{@{}lrrrl@{}}
\toprule
实现 & 参数距离 & 最大间隔漂移 & 越界样本 & 任务约束 \\
\midrule
O/B & 0.29904 & 5.36907 & 1/512 & 超过0.5 \\
A/B & 0.36414 & 39.92018 & 1/512 & 超过0.5 \\
\bottomrule
\end{tabular}
}

\newcommand{\NufiltTable}{
\begin{tabular}{@{}lrrr@{}}
\toprule
方法 & 无攻击 & \shortstack{AGA处理\\保护专家} & \shortstack{AGA处理\\全部专家} \\
\midrule
无保护 & 99.43 / 94.45 & 99.39 / 94.68 & 99.39 / 93.75 \\
偏置扰动 & 99.43 / 94.45 & 99.39 / 94.68 & 99.39 / 93.75 \\
补偿 A/B$^\dagger$ & 99.33 / 94.46 & 99.14 / 94.61 & 99.08 / 93.72 \\
\bottomrule
\end{tabular}
}

\newcommand{\EightTable}{
\begin{tabular}{@{}lrrr@{}}
\toprule
八专家TA0.2 & MNIST & DTD & SVHN \\
\midrule
无保护 & 93.23 & 52.34 & 73.90 \\
O/B$^\dagger$ & 54.22 & 50.80 & 44.63 \\
A/B$^\dagger$ & 59.16 & 50.59 & 51.97 \\
\bottomrule
\end{tabular}
}

\newcommand{\PhaseOneRecoveryTable}{
\begin{tabular}{@{}lrrr@{}}
\toprule
目标／伙伴 & TA+AGA & NUFILT+AGA & NUFILT无保护 \\
\midrule
MNIST / SVHN & 9.80 / 6.70 & 99.39 / 94.68 & 99.39 / 94.68 \\
DTD / MNIST & 2.13 / 21.09 & 70.59 / 99.45 & 70.59 / 99.45 \\
SVHN / MNIST & 6.70 / 9.80 & 94.40 / 99.25 & 94.40 / 99.25 \\
\bottomrule
\end{tabular}
}

\begin{document}

\begin{frame}[plain]
\vspace{12mm}
{\color{accent}\small Post-hoc Model Merge Protection\par}
\vspace{5mm}
{\LARGE\bfseries 模型融合的后处理防御\par}
\vspace{4mm}
{\large 两种方案与攻击评估\par}
\vspace{13mm}
{\color{muted} CLIP ViT-B/32 \quad RTX 3090\par}
{\color{muted}\small 2026年10月16日\par}
\end{frame}

\begin{frame}{研究问题}
\key{专家独立使用时保持能力，未经授权融合后降低被保护任务性能。}
\[
\theta_{\rm release}=T_c(\theta_e),\qquad
\theta_{\rm merge}=\theta_0+\tfrac12
\big[(\theta_{\rm release}-\theta_0)+(\theta_p-\theta_0)\big].
\]
\begin{itemize}\setlength{\itemsep}{6pt}
\item 已微调专家$\theta_e$冻结；保护发生在模型发布前。
\item 构造仅使用目标专家、公开基座及目标训练数据，不读取未来伙伴。
\item 同时评价\textbf{被保护任务与伙伴任务}，区分融合和攻击恢复。
\end{itemize}
\vspace{3mm}
两套方案：\textbf{投影核空间偏置扰动}；\textbf{任务子空间约束的跨模块补偿}。
\source{评价使用历史测试集，冻结防御后执行；不将其称为首次盲测。TA为Task Arithmetic。}
\end{frame}

\begin{frame}{方案一：投影核空间中的偏置扰动}
CLIP末端特征为$t=\operatorname{LN}_{\gamma,b}(h)P_e$，
$P_e\in\mathbb R^{768\times512}$。选择
\[
v^\top P_e=0,\qquad b'=b+\alpha v
\quad\Longrightarrow\quad t'=t.
\]
独立模型不读取该方向；融合后读出矩阵改变，扰动变为
$\delta t_{\rm merge}=\frac{\alpha}{2}v^\top P_{\rm merge}$。
\vspace{2mm}
\begin{center}\small
\PhaseOneTable
\end{center}
\small
AGA原版未修改关键LN偏置；这些结果不能证明抗任意恢复。
\source{准确率（\%）；AGA处理全部参与专家。该方案与第二方案的0.5距离预算分列，不能直接作同预算优劣比较。来源：results/phase3/results.json。}
\end{frame}

\begin{frame}{方案二：任务子空间约束下的跨模块补偿}
对任务矩阵$\Delta W=W_e-W_0=U\Sigma V^\top$构造奇异方向探针，
仅求解发布变换系数$c$。
\[
\max_c D(T_c(\theta_e),\theta_e)
\quad\text{s.t.}\quad
D\le0.5,\quad R_{\mathcal C\cup\mathcal P}(c)\le0.5.
\]
\begin{itemize}\setlength{\itemsep}{5pt}
\item $R$：自然训练输入与奇异探针上的最坏相对分类间隔变化。
\item 内层选择最坏任务输入／竞争类别；外层在有限可行域内增大参数距离，并使用公开基座代理诊断分类响应。
\item O/B联合修改Attention输出和MLP输出；A/B联合修改MLP输入和输出。
\end{itemize}
\key{原任务中的联合补偿，可能在融合后的表示与门控变化下失配。}
\source{SVD是结构工具，不是功能保持证明；有限内层与公开代理不构成未知伙伴下的安全保证。使用发布系数梯度，不能称零训练成本。}
\end{frame}

\begin{frame}{真实融合与AGA结果}
\small MNIST为被保护任务，SVHN为伙伴；每格为\textbf{MNIST / SVHN}准确率（\%）。
\vspace{3mm}
\begin{center}\footnotesize
\MainTable
\end{center}
\vspace{2mm}
\small
PaRaMS完整配置的直接融合破坏可被AGA恢复。
O/B、A/B的干扰仍然存在，但需要单独核验功能约束。
\source{TA系数0.5；两列AGA分别处理保护专家、全部专家。*不同预算参考；$^\dagger$已发现任务保持预算反例。PaRaMS为完整官方配置。}
\end{frame}

\begin{frame}{准确率保持，不等于最坏功能误差合格}
\begin{center}\small
\UtilityTable
\end{center}
\vspace{3mm}
\small
同一训练样本的教师7/9分数间隔仅为0.02034：
\begin{itemize}\setlength{\itemsep}{5pt}
\item O/B使该间隔增大，预测仍为7；相对漂移仍超过0.5。
\item A/B使间隔变为负值，预测由7变为9；同样违反约束。
\end{itemize}
\vspace{2mm}
\key{两者保留为真实效果参考，不能宣称全任务满足双0.5预算。}
\source{固定512张训练图的只读审计；完整模型重放误差为0。历史有限确认保留，不使用暴露见证修补权重。来源：utility.json、utility\_witnesses.json。}
\end{frame}

\begin{frame}{攻击方法：依据、权限与复现状态}
\scriptsize
\setlength{\tabcolsep}{4pt}
\renewcommand{\arraystretch}{1.28}
\begin{tabular}{@{}p{3.4cm}p{6.2cm}p{3.9cm}@{}}
\toprule
攻击／恢复流程 & 关键权限与操作 & 当前状态\\
\midrule
AGA & 公开基座；Attention求解＋MLP匹配 & 官方代码已复现\\
PaRaMS自适应恢复 & 基座与防御知识；恢复排列／尺度 & MLP官方对齐入口已核对；未运行\\
MergeLock配对对齐 & 伙伴参照；Q/K、V/O耦合正交对齐 & 原文核对；官方库仅README\\
MergeBarrier解码／层替换 & 基座、变换／修改位置知识 & 原文与附录核对；替换仅历史诊断\\
Unmask / GradErase & 参数减法／任务数据＋投影梯度再训练 & 正式论文核对；未复现\\
蒸馏、融合后SFT & 教师查询／任务数据与额外训练 & 未运行；强权限单列\\
ProDistill用于恢复 & 校准输入＋逐层特征蒸馏，学习融合系数 & 官方Vision代码已读；未运行\\
\bottomrule
\end{tabular}
\vspace{3mm}
\small
\textbf{TA、TIES、NUFILT是融合算法。}攻击后可以接不同融合算法，但不把它们计为新增恢复攻击。
\source{PaRaMS \S6；MergeLock \S4.2.2；MergeBarrier \S4.1、附录D；MergeGuard \S7；Trap² \S6.4。完整链接与固定源码见notes/phase3\_attack\_review.md。}
\end{frame}

\begin{frame}{更换融合算法后的保护边界}
\small
NUFILT原作者内核＋现有CLIP权重适配；保护MNIST，伙伴SVHN。
两个任务顺序均评价，下面展示MNIST先加入的固定顺序。
\vspace{2mm}
\begin{center}\small
\NufiltTable
\end{center}
\vspace{2mm}
\begin{itemize}\setlength{\itemsep}{4pt}
\item 其默认合并不读取关键LN偏置，方案一与对应无保护合并模型相同。
\item 两个顺序中，A/B距匹配无保护MNIST控制仅0.06--0.31个百分点。
\end{itemize}
\source{每格为MNIST / SVHN准确率（\%）。无图片参与融合门控求解；原专家和防御不更新。此处是融合方法评价，不是新增攻击；未复现完整FusionBench论文表。}
\end{frame}

\begin{frame}{当前结论与后续实验}
\begin{itemize}\setlength{\itemsep}{11pt}
\item \textbf{两种机制均观察到真实融合干扰。}\\
方案一有跨被保护任务结果；方案二仍主要限于MNIST。
\item \textbf{安全边界已经明确。}\\
方案一依赖融合参数覆盖；方案二存在任务预算反例，通用性与创新性尚未确立。
\item \textbf{下一步补攻击证据，不改冻结防御。}\\
优先补齐可核验的对齐／解码恢复；蒸馏和再训练按更强权限单独评价。
\end{itemize}
\source{全部既有权重、原始预测和失败证据保留。当前完成的是固定攻击与融合设置，未宣称覆盖所有攻击。}
\end{frame}
\end{document}
